\chapter{Small-Bowel Resection}

\section*{Introduction \& Clinical Indications}
Small-bowel resection removes a diseased or nonviable segment of jejunum or ileum and may be followed by an anastomosis or a stoma. Indications include ischemia or strangulation, perforation, obstruction, trauma, tumors, complicated diverticular or inflammatory disease, and strictures or fistulas that cannot be managed medically or endoscopically. The operation is urgent when there is threatened bowel, perforation, peritonitis, or sepsis; elective treatment is individualized in Crohn disease to preserve bowel length.

\section*{Relevant Surgical Anatomy}
The small bowel extends from the ligament of Treitz to the ileocecal junction. Its mesentery contains branches of the superior mesenteric artery and vein, lymphatics, and autonomic nerves. The surgeon assesses bowel viability, the mesenteric arcades, adjacent colon, ureters, and abdominal wall. Preserving adequate length and blood supply is essential to avoid short-bowel syndrome and anastomotic failure.

\section*{Preoperative Workup \& Preparation}
Assessment includes hemodynamics, sepsis, abdominal examination, nutritional status, comorbidities, medication and anticoagulant review, and imaging such as CT when time permits. Blood count, electrolytes, renal function, lactate, coagulation testing, and type and screen are selected according to urgency and risk. Resuscitation, antibiotics for contamination or suspected infection, venous-thromboembolism prevention, and nutritional optimization are planned. Consent includes anastomotic leak, abscess, ileus, obstruction, stoma, short-bowel syndrome, reoperation, and death.

\section*{Surgical Technique \& Procedure}
\begin{enumerate}
  \item The patient receives general anesthesia and is positioned supine. The abdomen is accessed by laparotomy or a minimally invasive approach according to the indication and stability.
  \item The bowel is examined from the ligament of Treitz to the ileocecal junction, and the diseased or nonviable segment and its cause are identified.
  \item Mesenteric vessels supplying the segment are divided with preservation of perfusion to the remaining bowel. The segment is removed with margins appropriate to the disease.
  \item If the patient is stable and the ends are healthy, a hand-sewn or stapled anastomosis is created. A stoma or delayed reconstruction may be safer with contamination, shock, poor perfusion, or high leak risk.
  \item Hemostasis, counts, bowel orientation, and the abdominal wall are checked before closure.
\end{enumerate}

\section*{Complications \& Risks}
Risks include bleeding, infection, wound complications, ileus, recurrent obstruction, anastomotic leak, intra-abdominal abscess, fistula, sepsis, venous thromboembolism, injury to adjacent structures, short-bowel syndrome, and recurrent disease. Crohn disease commonly recurs and requires ongoing medical and surgical follow-up.

\section*{Postoperative Care \& Recovery}
Monitoring includes vital signs, abdominal examination, urine output, pain, wound, stoma if present, bowel function, and laboratory trends. Oral intake is advanced according to clinical condition and tolerance; early oral feeding is appropriate for many patients when there is no contraindication. Mobilization, thromboprophylaxis, analgesia, and nutrition support are coordinated. Fever, tachycardia, worsening pain, distension, vomiting, feculent drainage, or failure to recover bowel function requires assessment for leak, abscess, or obstruction. Histopathology guides further treatment when tissue is removed.

\section*{Outcomes \& Prognosis}
Prognosis depends on the cause, urgency, bowel length and viability, contamination, nutrition, comorbidity, and anastomotic safety. Recovery is often good after uncomplicated elective resection, while emergency ischemia, perforation, malignancy, and Crohn disease carry distinct risks. Recurrence rates are disease-specific and cannot be represented by one universal percentage.

\section*{Additional Clinical Considerations}
The operative decision is guided by physiology as well as anatomy. A patient with shock, acidosis, rising lactate, peritonitis, or suspected bowel necrosis needs resuscitation and source control without unsafe delays for extensive testing. Broad-spectrum antibiotics, fluid or blood-product support, electrolyte correction, and early surgical involvement are coordinated. CT can identify obstruction, ischemia, perforation, abscess, and the likely transition point when the patient is stable enough for imaging, but a reassuring scan does not override worsening clinical findings.

Bowel viability is assessed using color, peristalsis, mesenteric pulsation, bleeding from the cut edge, and sometimes fluorescence or other adjuncts. These methods are imperfect, particularly after strangulation or reperfusion. The surgeon balances removal of nonviable bowel against preservation of length. Excessive resection can cause short-bowel syndrome, dependence on parenteral nutrition, dehydration, electrolyte loss, and impaired quality of life. A second-look operation may be considered when viability remains uncertain or the patient is too unstable for a definitive reconstruction.

An anastomosis is safer when the bowel ends are well perfused, tension-free, appropriately sized, and not exposed to uncontrolled contamination or severe physiologic compromise. A stoma or temporary abdominal closure may be safer in selected emergency cases. The decision is individualized and may be revisited after resuscitation. A leak can present with tachycardia, fever, increasing pain, ileus, leukocytosis, or subtle deterioration before obvious peritonitis; early imaging, drainage, antibiotics, and reoperation are considered according to the patient's condition.

Nutrition is part of treatment. Patients with obstruction, prolonged fasting, cancer, Crohn disease, or major sepsis are assessed for malnutrition and refeeding risk. Enteral nutrition is generally preferred when the gut is functional, while parenteral nutrition may be needed for prolonged ileus, severe short-bowel syndrome, or inability to meet requirements enterally. Stoma output, urine output, sodium, magnesium, renal function, and hydration require close monitoring after ileostomy creation, especially during the first weeks.

The pathology report may identify malignancy, ischemic injury, inflammatory disease, infection, or another cause and can change staging and adjuvant treatment. Crohn disease frequently requires long-term medical therapy even after technically successful resection; recurrence may occur at the anastomosis or elsewhere. Patients should receive individualized advice on diet progression, wound and stoma care, activity, medication, smoking cessation, and warning signs. Persistent vomiting, progressive distension, feculent drainage, fever, or reduced stoma output requires urgent assessment.

\section*{Additional Follow-Up Considerations}
After discharge, the patient should understand the expected bowel pattern, wound care, diet progression, and the distinction between normal postoperative fatigue and deterioration. A temporary ileus may delay intake, but worsening distension, persistent vomiting, inability to pass stool or gas, fever, tachycardia, or increasing abdominal pain requires evaluation. Patients with a new stoma need education about appliance fitting, output measurement, skin protection, dehydration, and when to contact the stoma team.

The anastomosis and the remaining bowel are affected by perfusion, nutrition, inflammation, medications, smoking, and systemic disease. Early mobilization and venous-thromboembolism prevention are balanced against the injury and any weight-bearing limitations. Antibiotics are continued only when indicated by contamination, infection, or the clinical course. Imaging and laboratory testing are guided by symptoms rather than performed on a fixed schedule for every uncomplicated recovery.

Crohn disease, ischemic disease, cancer, and adhesional obstruction require different follow-up plans. Crohn disease may recur despite a patent anastomosis and should be managed with gastroenterology and nutritional care. Cancer follow-up follows the stage and pathology, including margin and node status. Ischemia prompts assessment of the vascular or embolic cause when relevant. Patients with extensive resection require monitoring of micronutrients, bile-acid diarrhea, renal stones, bone health, and the possibility of intestinal rehabilitation.

\section*{Nutrition, Stoma, and Disease-Specific Follow-Up}
Nutritional assessment should continue after discharge. High-output ileostomy can cause sodium and water depletion even when the patient feels well, and patients may need oral rehydration solutions rather than plain water alone. Persistent diarrhea, steatorrhea, weight loss, edema, or fatigue can signal malabsorption or micronutrient deficiency. A dietitian may guide gradual diet advancement, oral supplements, vitamin replacement, and the timing of parenteral nutrition when intestinal failure is present.

The pathology and operative findings determine whether additional treatment is required. A malignant specimen may require staging scans, molecular testing, oncology referral, and systemic treatment. An ischemic specimen may prompt evaluation for embolic, thrombotic, low-flow, or medication-related causes. Crohn disease requires a plan for endoscopic or imaging surveillance and medical prevention of recurrence. Adhesive obstruction may recur after any abdominal operation, so recurrent cramping, vomiting, or distension warrants early assessment rather than repeated self-treatment.

\section*{Readmission and Recovery Risks}
Early postoperative deterioration can be subtle. Persistent tachycardia, a new oxygen requirement, falling urine output, increasing abdominal pain, or failure of inflammatory markers to improve may precede a clinically obvious leak or abscess. A low threshold for repeat examination and cross-sectional imaging is appropriate when the course is not progressing as expected. Treatment may include antibiotics, percutaneous drainage, nutritional support, endoscopic intervention, or reoperation.

The patient should receive individualized advice about lifting, driving, work, diet, and medication. Recovery after minimally invasive elective surgery may be shorter than after emergency laparotomy, but the underlying disease often determines the timeline. Adhesions can cause obstruction months or years later, while an anastomosis for cancer or Crohn disease needs disease-specific surveillance. A stoma reversal is considered only after healing, adequate sphincter and bowel function, and recovery from the original illness have been assessed.

\section*{Patient Education}
Patients should know whether the operation was an anastomosis or a stoma, how much bowel was removed, and which pathology results remain pending. They should seek urgent care for persistent vomiting, inability to keep fluids down, severe or worsening pain, fever, fainting, black or bloody stool, wound separation, or sudden reduction in stoma output. These instructions are particularly important after discharge from an emergency admission, when complications may develop after the initial improvement.

\section*{Multidisciplinary Care}
Emergency resection often requires coordination among surgery, anesthesia, critical care, radiology, infectious disease, nutrition, gastroenterology, oncology, and stoma specialists. The plan may change as perfusion and physiology improve. Clear handover records the anastomosis or stoma, remaining bowel length when clinically important, cultures, antibiotics, pathology, drains, and triggers for re-imaging or reoperation.

\section*{Documentation and Safety}
The record should state the indication, bowel viability, length and location of the resection, anastomosis or stoma, contamination, drains, cultures, and pathology plan. This information supports postoperative decisions about nutrition, hydration, antibiotics, surveillance, and the safety of later reconstruction. The patient should know which symptoms require urgent review and who is responsible for pathology and follow-up.

\section*{Practical Follow-Up}
Follow-up should include wound healing, bowel function, hydration, nutrition, and review of pathology. A stoma nurse or dietitian may be needed after ileostomy or extensive resection. The patient should have a clear plan for urgent assessment if pain, vomiting, fever, distension, or stoma-output changes develop.

\section*{Safety Summary}
The operation is urgent when bowel is threatened, but reconstruction is tailored to perfusion and physiology. Early recognition of leak, sepsis, obstruction, dehydration, and nutritional failure improves the chance of recovery.

\section*{Final Safety Note}
An uncomplicated recovery is not assumed when pain, tachycardia, fever, vomiting, distension, or stoma-output change occurs. Early surgical review helps identify leak, abscess, obstruction, sepsis, dehydration, and nutritional failure.

\section*{Completion Checklist}
Document the indication, bowel segment removed, anastomosis or stoma, contamination, pathology, nutrition plan, and warning symptoms. A clear handover reduces delays in recognizing leak, abscess, obstruction, dehydration, or recurrent disease.

\section*{Handoff}
The discharge handoff should identify the anastomosis or stoma, pathology status, nutrition needs, and the responsible surgical contact for deterioration.

\section*{Final Note}
Recovery and follow-up remain disease-specific, and deterioration warrants early surgical assessment.

\section*{References}
\begin{enumerate}
  \item American Society of Colon and Rectal Surgeons. Clinical Practice Guideline for the Surgical Management of Crohn's Disease.
  \item European Society for Clinical Nutrition and Metabolism. Practical Guideline: Clinical Nutrition in Surgery.
  \item Townsend CM, et al., editors. Sabiston Textbook of Surgery. Elsevier.
\end{enumerate}

\clearpage
